\documentclass[aspectratio=169,12pt]{beamer}
% Build from the repository root: make lecture-08-slides
\usepackage[T1]{fontenc}
\usepackage{lmodern}
\usepackage{amsmath,amssymb,booktabs,array}
\usepackage{graphicx}
\usefonttheme{professionalfonts}
\definecolor{Ink}{HTML}{16324F}
\definecolor{Blue}{HTML}{1F4E79}
\definecolor{Orange}{HTML}{C46A13}
\definecolor{Teal}{HTML}{007D80}
\setbeamercolor{normal text}{fg=Ink,bg=white}
\setbeamercolor{frametitle}{fg=Ink}
\setbeamercolor{structure}{fg=Blue}
\setbeamerfont{frametitle}{size=\large,series=\bfseries}
\setbeamersize{text margin left=8mm,text margin right=8mm}
\setbeamertemplate{navigation symbols}{}
\setbeamertemplate{headline}{}
\setbeamertemplate{footline}{}
\addtobeamertemplate{footline}{}{\vspace*{3mm}}
\setbeamertemplate{frametitle}{\vspace{3mm}\insertframetitle\par\vspace{1mm}}
\newcommand{\slidespacing}{%
  \setlength{\parskip}{0.2em}%
  \setlength{\abovedisplayskip}{0.3em}%
  \setlength{\belowdisplayskip}{0.3em}%
  \setlength{\abovedisplayshortskip}{0.1em}%
  \setlength{\belowdisplayshortskip}{0.2em}}
\newcommand{\Prob}{\mathbb P}
\newcommand{\Expect}{\mathbb E}
\DeclareMathOperator{\Var}{Var}
\DeclareMathOperator{\Cov}{Cov}
\newcommand{\KL}{D_{\mathrm{KL}}}
\newcommand{\CE}{\operatorname{CE}}
\newcommand{\plotcaption}[1]{\par\smallskip{\small #1\par}}
\setbeamertemplate{blocks}[default]
\setbeamerfont{block title}{size=\normalsize,series=\bfseries}
\setbeamerfont{block body}{size=\normalsize}
\newenvironment{coursebox}[2]{%
  \begingroup
  \setbeamercolor{block title}{fg=#1,bg=#1!13!white}
  \setbeamercolor{block body}{fg=Ink,bg=#1!4!white}
  \begin{block}{#2}
  \setlength{\parskip}{0.1em}
  \setlength{\abovedisplayskip}{0.3em}
  \setlength{\belowdisplayskip}{0.3em}
  \setlength{\abovedisplayshortskip}{0.1em}
  \setlength{\belowdisplayshortskip}{0.1em}
}{\end{block}\endgroup}
\newenvironment{definitionbox}[1]
  {\begin{coursebox}{Blue}{Definition: #1}}{\end{coursebox}}
\newenvironment{resultbox}[1]
  {\begin{coursebox}{Teal}{Result: #1}}{\end{coursebox}}
\newcommand{\optfigure}[1]{\par\begin{center}
  \includegraphics[width=.96\linewidth,height=4.45cm,keepaspectratio]{../../figures/ch8/#1.pdf}
  \end{center}}
\newcommand{\question}[1]{\par\medskip{\color{Blue}\bfseries #1\par}}
\AtBeginEnvironment{frame}{\slidespacing}
\title{Optimization for Machine Learning}
\hypersetup{pdftitle={Lecture 8: Optimization for Machine Learning},
  pdfsubject={Scientific Computing 2026},pdfauthor={}}
\begin{document}

{\setbeamercolor{background canvas}{bg=Ink}
\begin{frame}[plain]
\color{white}
{\large Scientific Computing 2026\par}
\vfill
{\LARGE\bfseries Optimization for Machine Learning\par}
\medskip
{\large Gradient descent, stochastic gradients, and Adam\par}
\vfill
{\large Lecture 8\par}\vspace{4mm}
\end{frame}}

\section{Fitting the homework model}
\begin{frame}{How expensive is the homework grid?}
In Homework 2 we tried 601 values of one parameter and kept the one
with the smallest loss.
\medskip

\begin{center}\renewcommand{\arraystretch}{1.35}
\begin{tabular}{@{}lr@{}}\toprule
Adjustable parameters & Loss evaluations\\\midrule
One & \(601\)\\
Two & \(601^2=361{,}201\)\\
Ten & \(601^{10}\approx6.15\times10^{27}\)\\\bottomrule
\end{tabular}\end{center}
\question{At the current parameter value, how can we tell which way to move?}
\end{frame}

\begin{frame}{One parameter controls both predictions}
The student predicts \(q_a(x)=\sigma(ax)\), where
\(\sigma(z)=1/(1+e^{-z})\). Start at \(a=\log 4\).
\begin{center}\renewcommand{\arraystretch}{1.4}
\begin{tabular}{@{}crr@{}}\toprule
Input \(x\) & Teacher probability \(p_x\) & Student probability \(q_a(x)\)\\\midrule
1 & \(0.80\) & \(4/5=0.80\)\\
2 & \(0.60\) & \(16/17\approx0.94\)\\\bottomrule
\end{tabular}\end{center}
\question{Would you increase or decrease \(a\)? Which prediction gets worse?}
\end{frame}

\begin{frame}{How much does a change in probability cost?}
For one input, write \(p\) for the teacher probability and \(q\) for the
student probability. Its loss is
\[
 \ell=-p\log q-(1-p)\log(1-q).
\]
\[
 \frac{\partial\ell}{\partial q}
   =-\frac p q+\frac{1-p}{1-q}
   =\frac{q-p}{q(1-q)} .
\]
When the student already matches the teacher, the slope is zero.
At \(x=2\) our student overshoots \(p\), so lowering \(q\) helps.
\end{frame}

\begin{frame}{The derivative with respect to the parameter}
The sigmoid has slope \(\sigma'(z)=\sigma(z)(1-\sigma(z))\).
Since \(q=\sigma(ax)\), we have \(dq/da=xq(1-q)\).
\begin{resultbox}{One example's gradient}
\[
 \ell'_x(a)=
 \underbrace{\frac{q-p_x}{q(1-q)}}_{\partial\ell/\partial q}
 \underbrace{xq(1-q)}_{dq/da}
 = x(q-p_x).
\]
\end{resultbox}
With equal input frequencies, \(L(a)=(\ell_1(a)+\ell_2(a))/2\).
The slope of the mean loss is the mean of the two slopes.
\end{frame}

\begin{frame}{Take one step by hand}
At \(a_0=\log4\), the two contributions are \(0\) and
\(2(16/17-0.60)\approx0.6824\). Thus \(L'(a_0)\approx0.3412\).
\begin{definitionbox}{Gradient descent}
Choose a learning rate \(\eta>0\). One update is
\[
 a_{k+1}=a_k-\eta L'(a_k).
\]
\end{definitionbox}
With \(\eta=1\): \(\quad a_1=1.3863-0.3412=1.0451.\)

The first-order change for a small step \(h\) is \(hL'(a)\). Choosing
\(h=-\eta L'(a)\) makes this quantity negative whenever \(L'(a)\ne0\).
\end{frame}

\begin{frame}{Each dot is one update}
\optfigure{student-descent}
The first step lowers the loss from \(0.8350\) to \(0.7315\) nats.
Further updates approach \(a_*\approx0.41962\). The homework grid
returned \(0.42086\).
\end{frame}

\section{Learning rate and curvature}
\begin{frame}{Can gradient descent fail on a parabola?}
Take \(F(w)=w^2/2\) and start at \(w_0=4\). Since the slope is \(w\),
\[
 w_{k+1}=w_k-\eta w_k.
\]
\question{Predict the next two parameter values for each learning rate.}
\[
  \eta=0.5 \qquad \eta=1.5 \qquad \eta=2.5
\]
Which runs cross zero? Which get closer to the minimum?
\end{frame}

\begin{frame}{The multiplier determines the next iterate}
For this parabola, every update multiplies \(w\) by \(1-\eta\).
\begin{center}\renewcommand{\arraystretch}{1.38}
\begin{tabular}{@{}rcll@{}}\toprule
\(\eta\) & Multiplier & First three values & Behaviour\\\midrule
0.5 & \(0.5\)  & \(4,\ 2,\ 1\)    & Converges\\
1   & \(0\)    & \(4,\ 0,\ 0\)    & One step\\
1.5 & \(-0.5\) & \(4,\ -2,\ 1\)   & Alternates and converges\\
2   & \(-1\)   & \(4,\ -4,\ 4\)   & Repeats\\
2.5 & \(-1.5\) & \(4,\ -6,\ 9\)   & Diverges\\\bottomrule
\end{tabular}\end{center}
\end{frame}

\begin{frame}{Opposite signs can give identical losses}
\optfigure{parabola}
With learning rates \(0.5\) and \(1.5\), both parameter magnitudes halve
on every step. Squaring removes the sign, so both losses fall by a factor
of four.
\end{frame}

\begin{frame}{The curvature sets the stable learning rates}
Let \(F(w)=cw^2/2\), with curvature \(F''(w)=c>0\).
\[
 w_{k+1}=(1-\eta c)w_k,\qquad w_k=(1-\eta c)^k w_0.
\]
\begin{resultbox}{Convergence on a scalar quadratic}
From any nonzero \(w_0\), the distance to zero shrinks exactly when
\[
 |1-\eta c|<1\quad\Longleftrightarrow\quad 0<\eta<2/c.
\]
\end{resultbox}
In Lecture 7, curvature measured how rapidly log-likelihood fell near a
peak. Here it tells us how quickly the slope changes during a step.
\end{frame}

\section{Two directions, one learning rate}
\begin{frame}{What if one direction is twenty times steeper?}
Let \(F(u,v)=(u^2+20v^2)/2\). Its gradient is the vector of partial
derivatives, \(\nabla F=(u,20v)\). Subtract \(\eta\nabla F\):
\[
 \begin{aligned}
    u_{k+1}&=(1-\eta)u_k,\\
    v_{k+1}&=(1-20\eta)v_k.
 \end{aligned}
\]
\question{Sketch two steps from \((4,1)\) using \(\eta=0.09\).}
The multipliers are \(0.91\) and \(-0.8\). Predict what happens to each
coordinate if the learning rate rises to \(0.11\).
\end{frame}

\begin{frame}{A change of 0.02 makes the path diverge}
\optfigure{valley-steps}
Left: \(v\) changes sign and retains \(80\%\) of its magnitude.
Right: \(v\) changes sign and grows by \(20\%\) each step.
Both panels show nine updates. Compare the vertical scales.
\end{frame}

\begin{frame}{Why the flat direction takes so long}
The Hessian collects the second partial derivatives. For this loss,
\[
 H=\begin{pmatrix}1&0\\0&20\end{pmatrix}.
\]
\begin{definitionbox}{Condition number of a positive-definite quadratic}
The condition number is the ratio of largest to smallest curvature
along the principal directions. Here \(\kappa=20/1=20\).
\end{definitionbox}
Stability along \(v\) forces \(0<\eta<0.1\). Along \(u\), every such step
retains more than \(90\%\) of the remaining distance. The flat direction
inherits the steep direction's speed limit.
\end{frame}

\section{Stochastic gradients and mini-batches}
\begin{frame}{At the best fit, the two examples still disagree}
Back to the student model. At \(a_*\approx0.41962\):
\begin{center}\renewcommand{\arraystretch}{1.45}
\begin{tabular}{@{}crrr@{}}\toprule
\(x\) & Teacher & Student & Gradient \(x(q-p_x)\)\\\midrule
1 & 0.800 & 0.6034 & \(-0.1966\)\\
2 & 0.600 & 0.6983 & \(+0.1966\)\\\bottomrule
\end{tabular}\end{center}
\question{One example is picked at random. Which way will \(a\) move?}
What would happen if the optimizer used both examples?
\end{frame}

\begin{frame}{A random step away from the optimum}
With \(\eta=0.35\), the next parameter is approximately
\[
 \begin{array}{lll}
 x=1: & 0.4196-0.35(-0.1966)=0.4884 & \text{higher},\\[.4em]
 x=2: & 0.4196-0.35(+0.1966)=0.3508 & \text{lower}.
 \end{array}
\]
Each outcome has probability \(1/2\). The expected update is zero.
\medskip

Both moves increase the \emph{full} loss. Its slope is zero at \(a_*\)
and its curvature is positive everywhere:
\[
 L''(a)=\tfrac12\sum_{x=1,2} x^2 q_a(x)\bigl(1-q_a(x)\bigr)>0.
\]
\end{frame}

\begin{frame}{One example gives an unbiased gradient}
Suppose \(L(w)=N^{-1}\sum_{i=1}^{N}\ell_i(w)\) is the mean loss on
\(N\) examples.
\begin{definitionbox}{Stochastic gradient descent (SGD)}
At the current \(w\), sample a uniform index \(I\) and take
\[
 w_{\rm next}=w-\eta\,\nabla\ell_I(w).
\]
\end{definitionbox}
At fixed \(w\), averaging over the index gives
\(\Expect[\nabla\ell_I(w)]=N^{-1}\sum_i\nabla\ell_i(w)=\nabla L(w)\).
One update needs one example gradient. A full update needs \(N\).
\end{frame}

\begin{frame}{How much quieter is a batch of 16?}
At a fixed parameter, let \(g_I\) be one component of a randomly sampled
example gradient. Average \(B\) independent draws with replacement:
\[
 g_B=\tfrac1B\sum_{j=1}^B g_{I_j}.
\]
\begin{resultbox}{Mini-batch variance}
\[
 \Expect[g_B]=\Expect[g_I],\qquad
 \Var(g_B)=\frac{B\Var(g_I)}{B^2}=\frac{\Var(g_I)}{B}.
\]
\end{resultbox}
At the homework optimum, the standard deviation falls from
\(0.1966\) for one draw to \(0.1966/4\approx0.0492\) for sixteen.
\end{frame}

\begin{frame}{Start all three runs at the optimum}
\optfigure{batch-paths}
Same learning rate, \(\eta=0.35\). The full gradient averages both distinct
inputs once. The blue run samples 16 inputs with replacement, so a batch
can contain unequal numbers of the two examples.
\end{frame}

\begin{frame}{What does an update cost?}
A training set contains one million examples.
\begin{center}\renewcommand{\arraystretch}{1.4}
\begin{tabular}{@{}lr@{}}\toprule
Update rule & Example gradients per update\\\midrule
One randomly chosen example & 1\\
A mini-batch of 100 & 100\\
Full gradient & 1,000,000\\\bottomrule
\end{tabular}\end{center}
\question{What should the horizontal axis show when comparing training cost?}
An update count measures progress per update. Examples processed or
elapsed time also account for the cost of each update.
\end{frame}

\section{Momentum and Adam}
\begin{frame}{Can past gradients help in the narrow valley?}
Along one coordinate, four gradients point the same way. Along the other,
they alternate. Keep half the old average and add half the new gradient.
Start at \(m_0=0\).
\begin{center}\renewcommand{\arraystretch}{1.35}
\begin{tabular}{@{}lrrrr@{}}\toprule
Update & 1 & 2 & 3 & 4\\\midrule
Persistent gradient & \(+1\)&\( +1\)&\( +1\)&\( +1\)\\
Its running average & .50& .75& .875& .9375\\[2pt]
Alternating gradient & \(+1\)&\( -1\)&\( +1\)&\( -1\)\\
Its running average & .50& \(-.25\)& .375& \(-.3125\)\\\bottomrule
\end{tabular}\end{center}
The persistent component gets a larger share of the next step.
\end{frame}

\begin{frame}{Momentum uses the running average for the step}
\begin{definitionbox}{Momentum}
Start with \(m_0=0\). At update \(k\ge1\), compute the current gradient
\(g_k\) at \(w_{k-1}\), then
\[
 m_k=\beta m_{k-1}+(1-\beta)g_k,\qquad
 w_k=w_{k-1}-\eta m_k.
\]
\end{definitionbox}
The coefficient \(0\leq\beta<1\) controls how much of the previous
average survives. Setting \(\beta=0\) recovers the ordinary gradient step.
\medskip

When the slope reverses, \(m_k\) can still point in the previous direction.
That is one way momentum can overshoot a minimum.
\end{frame}

\begin{frame}{Twelve steps on the same quadratic}
\optfigure{momentum}
Both paths start at \((4,1)\). Left: \(\eta=0.09\).
Right: \(\eta=0.16,\ \beta=0.5\). The larger dot marks step 12.
These paths illustrate the stated settings.
\end{frame}

\begin{frame}{Adam keeps a scale for each parameter}
Consider two parameters with first gradients \(10\) and \(0.1\).
An ordinary gradient step moves one parameter 100 times as far.
\medskip

For each parameter, Adam tracks two quantities separately:
\[
 \begin{aligned}
 m_k&=\beta_1 m_{k-1}+(1-\beta_1)g_k
      &&\text{average gradient},\\
 s_k&=\beta_2 s_{k-1}+(1-\beta_2)g_k^2
      &&\text{average squared gradient}.
 \end{aligned}
\]
Both start at zero, with \(0\le\beta_1,\beta_2<1\).
Even when signs alternate, the squared gradients accumulate.
\end{frame}

\begin{frame}{Adam's first step by hand}
\begin{definitionbox}{Adam update for one coordinate}
Correct the two averages for their zero initialization and take a step:
\[
 \widehat m_k=\frac{m_k}{1-\beta_1^k},\quad
 \widehat s_k=\frac{s_k}{1-\beta_2^k},\quad
 w_k=w_{k-1}-\eta\frac{\widehat m_k}{\sqrt{\widehat s_k}+\varepsilon}.
\]
\end{definitionbox}
At \(k=1\), the corrected values are \(g_1\) and \(g_1^2\).
The two gradients \(10\) and \(0.1\) give ratios
\[
  10/(10+\varepsilon)\quad\text{and}\quad 0.1/(0.1+\varepsilon).
\]
For small positive \(\varepsilon\), both updates have size close to \(\eta\).
\end{frame}

\begin{frame}{Which change would you try?}
\textbf{1.} The student model starts at its best fit. A single-example
update raises the full loss. What gradient would you inspect to check
whether the implementation is correct?
\bigskip

\textbf{2.} A full-gradient run on \(F(u,v)=(u^2+20v^2)/2\)
diverges with \(\eta=0.11\). Someone suggests increasing batch size.
What would you change first, and why?
\end{frame}

\begin{frame}{Two patterns in a training run}
\textbf{At the homework optimum}, inspect both example gradients.
Expect approximately \(-0.1966\) and \(+0.1966\). Their average vanishes.
A step based on either one moves away from the minimum.
\bigskip

\textbf{On the quadratic}, inspect the steep coordinate.
With \(\eta=0.11\), its multiplier is \(-1.2\). The run already uses the
full gradient. Reducing the learning rate to \(0.09\) brings the
multiplier to \(-0.8\) and the oscillations decay.
\end{frame}

\end{document}
